\section{Chronological partial source expansions}

We follow the chronological expansion in
\cite[proof of Theorem~5.2]{OpenAI2026PolynomialPEPS},
\texttt{04-compression.tex}, lines 342--434. Each gate has its own
participating parties and the common source data constructed above.
The following statements establish the exact partial expansion and
its scalar cost. The estimate for corrected source operators requires
the subsequent separation and second-moment arguments.

\begin{definition}[A chronological composition]
    \label{def:peps_source_circuit}
    \lean{TNLean.PEPS.PairEffect.SourceCircuit}
    \lean{TNLean.PEPS.PairEffect.SourceCircuit.eval}
    \lean{TNLean.PEPS.PairEffect.SourceCircuit.IsAllowed}
    \leanok
    \uses{def:peps_prepared_source_gate,def:peps_party_registers}
    Let $P$ be any set of parties. A chronological composition from
    register list $a$ to register list $b$ is formed from identities,
    consecutive composition, private operations, exchanges of adjacent
    registers, and nonprivate gates. A gate has a finite participant
    set $C$, an injection $C\hookrightarrow P$, and common finite
    source data on $C$. It acts on its prescribed input registers and
    as the identity on all remaining registers. Any composition may
    also be extended by an idle register. A composition is allowed
    when each private operation is a contraction. Its complete
    operator $W:\mathcal H_a\to\mathcal H_b$ uses the aggregate
    operator of each nonprivate gate.
\end{definition}

\begin{theorem}[Norm of the complete composition]
    \label{thm:peps_source_circuit_norm}
    \lean{TNLean.PEPS.PairEffect.SourceCircuit.norm_eval_le_one}
    \leanok
    \uses{def:peps_source_circuit}
    Every allowed chronological composition satisfies $\|W\|\leq1$.
\end{theorem}

\begin{proof}
    \leanok
    \uses{thm:peps_prepared_gate_grouped}
    Each gate is a contraction. Consecutive composition, extension
    by an identity, and unitary identifications of the register order
    preserve this bound. Induction on the composition proves the claim.
\end{proof}

\begin{definition}[Gate occurrences and source positions]
    \label{def:peps_source_circuit_positions}
    \lean{TNLean.PEPS.PairEffect.SourceCircuit.gateLocations}
    \lean{TNLean.PEPS.PairEffect.SourceCircuit.gateLocationsFintype}
    \lean{TNLean.PEPS.PairEffect.SourceCircuit.slotCount}
    \lean{TNLean.PEPS.PairEffect.SourceCircuit.sourceLocations}
    \lean{TNLean.PEPS.PairEffect.SourceCircuit.sourceLocationsFintype}
    \lean{TNLean.PEPS.PairEffect.SourceCircuit.participants}
    \lean{TNLean.PEPS.PairEffect.SourceCircuit.endpoints}
    \lean{TNLean.PEPS.PairEffect.SourceCircuit.branchLabels}
    \lean{TNLean.PEPS.PairEffect.SourceCircuit.branchLabelsFintype}
    \lean{TNLean.PEPS.PairEffect.SourceCircuit.branchCoefficient}
    \leanok
    \uses{def:peps_source_circuit}
    Let $\mathcal G$ be the finite set of occurrences of nonprivate
    gates. Repeated uses of the same gate give distinct elements of
    $\mathcal G$. For $g\in\mathcal G$, let $C_g\subseteq P$ be its
    participant set, $E_g$ its ordered source slots, $\Xi_g$ its
    original branch set, and $c_{g,\xi}$ its original coefficients.
    The source positions are
    $\mathcal E=\{(g,e):g\in\mathcal G,\ e\in E_g\}$.
    Each position has the two original endpoints of its slot,
    transported along the participant injection. No monomial label
    enters the definition of $\mathcal E$.
\end{definition}

\begin{theorem}[Incidence of the original source positions]
    \label{thm:peps_source_circuit_incidence}
    \lean{TNLean.PEPS.PairEffect.SourceCircuit.endpoints_mem_participants}
    \lean{TNLean.PEPS.PairEffect.SourceCircuit.endpoints_ne}
    \leanok
    \uses{def:peps_source_circuit_positions}
    The two endpoints of $(g,e)$ are distinct and both belong to $C_g$.
\end{theorem}

\begin{proof}
    \leanok
    \uses{def:peps_prepared_source_gate}
    Each original slot has distinct endpoints in its gate's participant
    set. The participant injection preserves distinctness. Consecutive
    composition retains the occurrence of the gate containing the slot.
\end{proof}

\begin{definition}[The partial branch choices]
    \label{def:peps_partial_source_choices}
    \lean{TNLean.PEPS.PairEffect.SourceCircuit.IsTouched}
    \lean{TNLean.PEPS.PairEffect.SourceCircuit.isTouched_gate_iff}
    \lean{TNLean.PEPS.PairEffect.SourceCircuit.Choices}
    \lean{TNLean.PEPS.PairEffect.SourceCircuit.choicesFintype}
    \lean{TNLean.PEPS.PairEffect.SourceCircuit.partialWord}
    \lean{TNLean.PEPS.PairEffect.SourceCircuit.coefficient}
    \lean{TNLean.PEPS.PairEffect.SourceCircuit.expandedEval}
    \leanok
    \uses{def:peps_source_circuit_positions,def:peps_affected_owner}
    Fix affected parties $\mathcal A\subseteq P$. A gate occurrence
    is touched when $C_g\cap\mathcal A\neq\varnothing$.
    Choose one label $\xi_g\in\Xi_g$ at each touched occurrence and
    no label at the other occurrences. Let $\Xi_{\mathcal A}$ be this
    finite set of choices. In the resulting composition $W_\xi$,
    touched gates are replaced by their chosen prepared branches,
    while untouched gates retain their complete operators. Identify
    all exterior owners as in Definition~\ref{def:peps_affected_owner}.
    The coefficient of the partial branch is
    $c_\xi=\prod_{g:\,C_g\cap\mathcal A\neq\varnothing}c_{g,\xi_g}$,
    and its weighted sum is $\sum_{\xi\in\Xi_{\mathcal A}}c_\xi W_\xi$.
    An untouched gate has one partial choice with coefficient one,
    even if its original branch set is empty.
\end{definition}

\begin{theorem}[Allowed partial branches and their actual sources]
    \label{thm:peps_partial_source_inventory}
    \lean{TNLean.PEPS.PairEffect.SourceCircuit.isAllowed_partialWord}
    \lean{TNLean.PEPS.PairEffect.SourceCircuit.expandedSources}
    \lean{TNLean.PEPS.PairEffect.SourceCircuit.sources_partialWord}
    \lean{TNLean.PEPS.PairEffect.SourceCircuit.mem_sources_partialWord}
    \leanok
    \uses{def:peps_partial_source_choices}
    If the original composition is allowed, each partial branch is
    allowed. Its source list is obtained from the chosen sources of
    the touched gates by retaining exactly those with at least one
    endpoint in $\mathcal A$, and then grouping the exterior owners.
    Each retained source has its original endpoint spaces and vector.
    The order and multiplicity of the source occurrences are preserved.
\end{theorem}

\begin{proof}
    \leanok
    \uses{thm:peps_affected_source_occurrences,thm:peps_prepared_gate_properties,thm:peps_grouped_block_map,thm:peps_prepared_gate_grouped}
    A chosen prepared branch is allowed. An untouched gate is a single
    contraction on exterior memory. The owner identification turns a
    pair with two exterior endpoints into a local normalized
    preparation. Induction on the chronological composition gives the
    asserted source list and allowedness.
\end{proof}

\begin{theorem}[Exact chronological operator identity]
    \label{thm:peps_partial_source_evaluation}
    \lean{TNLean.PEPS.PairEffect.Word.eval_sum_appendTail}
    \lean{TNLean.PEPS.PairEffect.SourceCircuit.eval_partial_gate}
    \lean{TNLean.PEPS.PairEffect.SourceCircuit.expandedEval_comp}
    \lean{TNLean.PEPS.PairEffect.SourceCircuit.expandedEval_frame}
    \lean{TNLean.PEPS.PairEffect.SourceCircuit.expandedEval_mapOwner}
    \leanok
    \uses{def:peps_partial_source_choices}
    Let $I_a,I_b$ be the canonical isometries identifying the input
    and output memories with their grouped-owner memories. Then
    \begin{align}
        \left(\sum_{\xi\in\Xi_{\mathcal A}}c_\xi W_\xi\right)I_a
         & = I_bW.
        \label{eq:peps_partial_source_evaluation}
    \end{align}
    The private operations need not be contractions for this identity.
\end{theorem}

\begin{proof}
    \leanok
    \uses{thm:peps_grouped_gate_sum,thm:peps_prepared_gate_grouped,thm:peps_grouped_block_map,def:peps_source_circuit}
    At a touched gate, sum its original branch expansion and extend
    the sum by the spectator identity. At an untouched gate, the
    single partial branch is its aggregate operator. The resulting
    identity is compatible with consecutive composition, since the
    two independent finite sums distribute and their coefficients
    multiply. It is also compatible with adjoining idle registers.
    Induction gives (\ref{eq:peps_partial_source_evaluation}).
\end{proof}

\begin{theorem}[Local contributions to the scalar cost]
    \label{thm:peps_partial_source_local_cost}
    \lean{TNLean.PEPS.PairEffect.SourceCircuit.sum_norm_coefficient_gate_of_touches}
    \lean{TNLean.PEPS.PairEffect.SourceCircuit.sum_norm_coefficient_gate_of_exterior}
    \lean{TNLean.PEPS.PairEffect.SourceCircuit.sum_norm_coefficient_comp}
    \lean{TNLean.PEPS.PairEffect.SourceCircuit.sum_norm_coefficient_frame}
    \lean{TNLean.PEPS.PairEffect.SourceCircuit.sum_norm_coefficient_id}
    \lean{TNLean.PEPS.PairEffect.SourceCircuit.sum_norm_coefficient_localMap}
    \lean{TNLean.PEPS.PairEffect.SourceCircuit.sum_norm_coefficient_swap}
    \leanok
    \uses{def:peps_partial_source_choices}
    A touched gate contributes $\sum_{\xi\in\Xi_g}|c_{g,\xi}|$ to
    the sum of absolute partial coefficients. An untouched gate,
    identity, private operation, or register exchange contributes one.
    Consecutive compositions multiply these sums. Adjoining an idle
    register leaves the sum unchanged.
\end{theorem}

\begin{proof}
    \leanok
    \uses{def:peps_partial_source_choices}
    Use the defining choice set at each elementary operation. For
    consecutive compositions, apply $|cd|=|c|\,|d|$ and distribute the
    sum over the Cartesian product of the two choice sets.
\end{proof}

\begin{theorem}[The exact product of the gate costs]
    \label{thm:peps_partial_source_product_cost}
    \lean{TNLean.PEPS.PairEffect.SourceCircuit.sum_norm_coefficient_eq_prod}
    \lean{TNLean.PEPS.PairEffect.SourceCircuit.sum_norm_coefficient_le_pow_card}
    \lean{TNLean.PEPS.PairEffect.SourceCircuit.sum_norm_coefficient_mul_conj_eq_prod_sq}
    \lean{TNLean.PEPS.PairEffect.SourceCircuit.sum_norm_coefficient_mul_conj_le_pow_card}
    \leanok
    \uses{def:peps_partial_source_choices}
    Let $\mathcal G_{\mathcal A}$ denote the touched gate occurrences
    and put $t=|\mathcal G_{\mathcal A}|$. Then
    \begin{align}
        \sum_{\xi\in\Xi_{\mathcal A}}|c_\xi|
         & =\prod_{g\in\mathcal G_{\mathcal A}}
        \sum_{\eta\in\Xi_g}|c_{g,\eta}|,
        \label{eq:peps_partial_source_product_cost}
    \end{align}
    and the sum of $|c_\xi\overline{c_\zeta}|$ over all independent
    ket and bra choices is the square of this product. If each touched
    gate's coefficient sum is at most $B$, these two costs are bounded
    by $B^t$ and $B^{2t}$, respectively. Every occurrence contributes
    separately, including repeated uses of one gate.
\end{theorem}

\begin{proof}
    \leanok
    \uses{thm:peps_partial_source_local_cost,def:peps_partial_source_choices}
    Induction on the chronological composition gives the product
    formula from the local contributions. Bound each nonnegative
    factor by $B$. The ket and bra sums are independent, and complex
    conjugation preserves absolute value, so their total cost is the
    square of the first sum.
\end{proof}

\begin{theorem}[The density expansion in orthonormal coordinates]
    \label{thm:peps_partial_source_density}
    \lean{TNLean.PEPS.PairEffect.SourceCircuit.expandedEval_eq_mapOwner}
    \lean{TNLean.PEPS.PairEffect.SourceCircuit.toMatrix_eval_eq_sum_partialWord}
    \lean{TNLean.PEPS.PairEffect.SourceCircuit.density_eval_eq_sum_partialWord}
    \leanok
    \uses{thm:peps_partial_source_evaluation}
    Choose finite orthonormal bases of the grouped input and output
    memories. Let $M$ be the matrix of $I_bWI_a^{-1}$ and $K_\xi$
    the matrix of the partial branch $W_\xi$. For every input matrix
    $\rho$,
    \begin{align}
        M & =\sum_{\xi\in\Xi_{\mathcal A}}c_\xi K_\xi,
        \label{eq:peps_partial_source_matrix}          \\
        M\rho M^\dagger
          & =\sum_{\xi,\zeta\in\Xi_{\mathcal A}}
        c_\xi\overline{c_\zeta}\,
        K_\xi\rho K_\zeta^\dagger.
        \label{eq:peps_partial_source_density}
    \end{align}
    In each summand, the untouched gates retain their complete
    operators. The matrix $\rho$ need not be positive or normalized.
\end{theorem}

\begin{proof}
    \leanok
    \uses{thm:peps_partial_source_evaluation}
    Compose (\ref{eq:peps_partial_source_evaluation}) with $I_a^{-1}$
    and take matrices. Distribute the finite sums in $M\rho M^\dagger$.
    Taking the adjoint conjugates each bra coefficient.
\end{proof}

\begin{theorem}[Fixed positions for every partial source inventory]
    \label{thm:peps_partial_source_order}
    \lean{TNLean.PEPS.PairEffect.SourceCircuit.sourceOrder}
    \lean{TNLean.PEPS.PairEffect.SourceCircuit.mem_sourceOrder}
    \lean{TNLean.PEPS.PairEffect.SourceCircuit.nodup_sourceOrder}
    \lean{TNLean.PEPS.PairEffect.SourceCircuit.sourceOrder_gate}
    \lean{TNLean.PEPS.PairEffect.SourceCircuit.sourceDims}
    \lean{TNLean.PEPS.PairEffect.SourceCircuit.sourceAt}
    \lean{TNLean.PEPS.PairEffect.SourceCircuit.sourceAt_isSome_iff}
    \lean{TNLean.PEPS.PairEffect.SourceCircuit.sourceAt_eq_some_spec}
    \lean{TNLean.PEPS.PairEffect.SourceCircuit.filterMap_sourceAt_comp}
    \lean{TNLean.PEPS.PairEffect.SourceCircuit.filterMap_sourceAt_gate}
    \lean{TNLean.PEPS.PairEffect.SourceCircuit.sources_partialWord_eq_filterMap}
    \lean{TNLean.PEPS.PairEffect.SourceCircuit.layout_sources_partialWord}
    \leanok
    \uses{def:peps_source_circuit_positions,def:peps_partial_source_choices}
    The positions $\mathcal E$ admit a fixed exhaustive ordered list
    without repetitions, in the preparation order of the actual
    composition. Consecutive compositions list the later composition's
    positions first; the original slot order within a gate is retained.
    Each position $e$ has fixed halfspaces $\C^{u_e}$ and $\C^{v_e}$.

    For every partial choice $\xi$, its source inventory is obtained
    from this fixed list by retaining exactly those positions with an
    endpoint in $\mathcal A$. A retained position has the grouped
    original endpoints, the fixed halfspaces, and its normalized
    chosen branch vector. Consequently the ordered layout of all
    surviving source registers is independent of $\xi$.
\end{theorem}

\begin{proof}
    \leanok
    \uses{thm:peps_partial_source_inventory,thm:peps_affected_source_occurrences}
    Induct on the chronological composition. The two occurrence sets
    of consecutive compositions are disjoint, so concatenating their
    lists preserves uniqueness. At a gate, list its prescribed slots.
    The participant injection preserves their distinct endpoints, and
    grouping owners removes a slot exactly when both endpoints are
    exterior. The remaining fields are its prescribed spaces and
    vector. Removing vectors from these records leaves a layout
    determined solely by the original positions and $\mathcal A$.
\end{proof}
